\subsection{Full-character reciprocity and a further two-moment obstruction}
\label{sec:r2603-center-moments}

The moment method also applies without conductor saturation. Let $K$ be
a quadratic dimension field and $t_K$ normalized relative trace from a
cyclotomic field containing all twists. The anomaly-free center Galois
identity is $\theta_{\pi_\sigma A}=\sigma^2(\theta_A)$. Since
$\sigma^2$ fixes $K$ and cyclotomic Galois automorphisms commute,
\begin{equation}\label{eq:quadratic-relative-moment-pairing}
 t_K(\theta_{\pi_\sigma A}^n)
   =\sigma^2\bigl(t_K(\theta_A^n)\bigr)=t_K(\theta_A^n)
 \qquad(n\geq1).
\end{equation}
Thus full-row Galois partners have equal relative moments even when
their twists differ. No exponent bound is needed.

\begin{theorem}\label{thm:r2603-center-moments}
The fusion ring $R7\_4\_2603$ admits no complex fusion categorification.
\end{theorem}
\begin{proof}
Proposition~\ref{prop:four-dihedral}, the two-dimension Galois theorem
\cite[Lemma~3.1]{Schopieray}, and Morita invariance of global dimension
reduce the problem to the pseudo-unitary case. In source order write
the basis as $(1,Y,Z,g,a,b,X)$. Put $s=\sqrt5$ and $u=\sqrt{-2}$.
Then
\[
 d=(1,4+2s,4+2s,1,1,1,2),\qquad D=80+32s,\qquad Y^*=Z.
\]
The invertibles form $V_4$. The complete character table is
\[
\begin{array}{c|rrrrrrr|c}
 &1&Y&Z&g&a&b&X&f_E\\\hline
 P&1&0&0&1&1&1&-2&8\\
 U&1&0&0&-1&1&-1&0&4\\
 V&1&0&0&-1&-1&1&0&4\\
 C_+&1&u&-u&1&-1&-1&0&8\\
 C_-&1&-u&u&1&-1&-1&0&8\\
 \overline d&1&4-2s&4-2s&1&1&1&2&80-32s\\
 d&1&4+2s&4+2s&1&1&1&2&80+32s
\end{array}
\]
All products and duality-defined codegrees in this table are checked
against the source tensor. Disjointness of $\mathbb Q(s)$ and
$\mathbb Q(u)$ gives an automorphism flipping $s$ and fixing every
other character. If $v\cdot d=A(v)+B(v)s$, its signed full-row action is
\[
 T(v)=v+(0,2,2,0,0,0,0)\frac{2A(v)-5B(v)}4.
\]
This integral involution follows from
\eqref{eq:quadratic-row-galois}; evaluation on every displayed
character verifies it directly. In particular the Galois image $J$
of the central unit has forgetful object $1+Y+Z$, of dimension $9+4s$.

\medskip\noindent\emph{Two possible induced-unit blocks.}
The other five nonunit constituents of $I(1)$ have dimensions
$10+4s$ three times and $20+8s$ twice. The unit Gram row is
$(7,8,8,3,1,1,2)$. A codegree-eight row must therefore be
\[
          1+h+pY+(2-p)Z,\qquad 0\leq p\leq2,
\]
with $h$ a nonidentity invertible. Each codegree-four row has at
most one $X$ summand; together they must supply the unit Gram
coefficient two. Hence each has form
\[
          1+k+X+qY+(4-q)Z,\qquad0\leq q\leq4.
\]
The five entries $h,k$ have multiset $\{g,g,g,a,b\}$.
Symmetry of the center $S$-matrix gives
\begin{equation}\label{eq:r2603-character-reciprocity}
 f_F\chi_E(F(A_F))=f_E\chi_F(F(A_E)).
\end{equation}
The $P$-value on a codegree-four row is zero. Reciprocity with
$U,V$ therefore forces the invertible in the $P$ row to be $g$.
Reciprocity with $C_+$ then forces $p_P=1$.

For $U$ or $V$ against $C_+$ or $C_-$, one side of
\eqref{eq:r2603-character-reciprocity} belongs to $\{0,16\}$;
the other has real part in $\{0,8\}$ and imaginary part proportional
to $(2q-4)u$. Both real parts must vanish and $q=2$. Thus the two
codegree-four rows use $a,b$, once each, and both $C$ rows use $g$.
Finally reciprocity between $C_+$ and $C_-$ gives $p_++p_-=2$.
Up to permuting equal-codegree rows, the complete nonunit block is
therefore
\begin{gather*}
 1+Y+Z;\\
 1+g+2Z,\quad1+g+Y+Z,\quad1+g+2Y;\\
 1+a+X+2Y+2Z,\quad1+b+X+2Y+2Z,
\end{gather*}
or the same block with all three middle rows equal to $1+g+Y+Z$.
This theoretical enumeration is independently checked by the verifier:
of fifty bounded Gram-compatible blocks, exactly these two admit a
complete codegree-respecting character assignment.

\medskip\noindent\emph{All root orders in a finite moment hull.}
Set $u_\pm=(-1\pm s)/4$. Every pair
$(t_K(\theta),t_K(\theta^2))$ lies in the convex hull of the five
rational vertices of Lemma~\ref{lem:quadratic-two-moment-hull}, the six
pairs
\[
 (u_+,u_-),\ (u_-,u_+),\ (-u_+,u_-),\ (-u_-,u_+),\
 (0,-u_+),\ (0,-u_-),
\]
and negative one-half times each of those six pairs.
Indeed, write the root order as $2^\alpha5^\beta m$, with $m$ odd
and prime to five. If $\beta=0$, ordinary Ramanujan averages lie in
the rational hull. If $\beta\geq2$, $m$ is not squarefree, or
$\alpha\geq3$, both moments vanish. The other cases have
$\beta=1$, $\alpha=0,1,2$. Their primitive-order $5,10,20$ pairs
are the six displayed points, multiplied by $\mu(m)/\varphi(m)$,
with a possible interchange of the signs. For $m>1$,
$\varphi(m)\geq2$, so this scale lies in $[-1/2,1]$.
This proves the seventeen-vertex hull for arbitrary root orders.

\medskip\noindent\emph{One universal integer inequality.}
Every residual central dimension is $z=A+Bs$, and $D/z$ must be
integral. The full ring of integers here is
$\mathbb Z[(1+s)/2]$. With $n=A^2-5B^2$, the exact condition is that
\[
 \frac{2(80A-160B)}n,\qquad\frac{2(32A-80B)}n
\]
are integers of the same parity. The Gram bounds and this arithmetic
condition give $465$ global row candidates. The exact rank-one PSD
test on each residual matrix, followed by full-row Galois pairing,
leaves $72$ and $76$ paired row types.

The second indicators are $(1,0,0,1,1,1,1)$. To see the positive
sign of $X$, the unit occurs once in $Y^*XY$, and its palindromic
form has sign $\nu_2(X)$. Then $X^2$ contains each invertible once,
so all their signs are positive. The dual pair $Y,Z$ has indicator
zero. As before, subtract the twist-one low rows from the Gram and
two induction trace identities. Use the seventeen-vertex hull and
the paired relative moments in
\eqref{eq:quadratic-relative-moment-pairing}. In the same
$21+12+12$ coordinate order as \eqref{eq:quadratic-moment-column},
multiply columns and targets by eight, rather than four, to clear
the hull denominators.

The following covector is written in its Gram, first-moment and
second-moment blocks:
\begin{align*}
 y_G={}&(38,58,26,63,76,8,57,69,-70,-72,7,19,-20,-100,\\[-2pt]
       &\hspace{10mm}-22,-100,-76,-100,-100,0,103),\\
 y_1={}&(-17,50,-17,50,-43,100,37,-100,37,-100,33,-100),\\
 y_2={}&(51,-100,51,-100,-72,100,-54,100,-54,100,-58,100).
\end{align*}
Every allowed column has dot product at least $12$ with $y$.
There are $1224$ and $1292$ columns in the two branches and $1292$
distinct columns in their union. Their prescribed targets have dot
products $-1336$ and $-1240$. Neither can be a nonnegative combination
of allowed columns. This contradicts the existence of the center.

The standard-library verifier
\texttt{verify\_r2603\_center\_moments.py} reconstructs the original
census tensor, checks its immutable hash, every character and its
full Galois transport, the complete low-block coverage, maximal-order
integrality, all residual columns, and the exact inequality above.
It uses the already audited exact Gram helpers and no optimizer.
Independent ordinary and optimized replays confirm every step.
The initial Galois reduction proves the claimed general-complex
exclusion.
\end{proof}
